\section{Conclusion}

We studied LoRA adapter compression as a post-training security intervention for backdoored language models.
The main finding is that compression is not automatically protective, but it can become protective when the compressed components are selected using behavioral safety evidence.
\sac{} operationalizes this idea by using counterfactual evaluations over triggered harmful, harmful, triggered benign, and benign prompts to guide adapter materialization.

On Qwen3.5-27B, \sac{} reduces triggered harmful ASR from 0.953 to 0.169, and a same-gate prune-then-INT8 variant reaches 0.172 while preserving MMLU at 0.822.
The selected adapter also transfers to AdvBench and HarmBench.
Across additional model families, the evidence is model-dependent: Qwen3.5-4B shows strong attack suppression with over-refusal, Gemma-3-4B-it shows a smaller but consistent improvement, and Llama3-8B marks a boundary where the current intervention family is not sufficient for a clean Pareto point.

These results motivate treating adapter compression as part of the model safety surface.
Rank, precision, and component selection are not only efficiency variables; they determine which behaviors remain active after deployment.
Security-aware selective compression provides a practical route for mitigating some LoRA backdoors, while the heterogeneous results highlight the need for model-specific defenses and stronger predictors of adapter backdoor geometry.
